\begin{textsample}{POS Dim 1 – pi – Score 53.00 – pi/pi\_em\_37.txt}  \label{ex:f1_pos_059}
Língua \textbf{Espanhola} Aprender a língua \textbf{espanhola} propicia a criação de novas formas de engajamento e participação dos \textbf{alunos} em um mundo social cada vez mais \textbf{globalizado} e plural, em que as fronteiras entre países e interesses pessoais, locais, regionais, nacionais e transnacionais estão cada vez mais difusas e contraditórias.
Assim, o estudo da Língua \textbf{Espanhola} possibilita aos \textbf{alunos} ampliar \textbf{horizontes} de comunicação e de intercâmbio cultural, científico e acadêmico e, nesse sentido, abre novos percursos de acesso, construção de conhecimentos e participação social.
É esse caráter formativo que inscreve a aprendizagem de \textbf{espanhol} em uma perspectiva de educação linguística, consciente e \textbf{crítica}, na qual as dimensões pedagógica e política são intrinsecamente ligadas.
Ensinar \textbf{espanhol} com essa finalidade tem, para o currículo, duas implicações importantes.
A primeira é que ela obriga a rever as relações entre língua, território e cultura, na medida em que os falantes de \textbf{espanhol} já não se encontram apenas nos países em que ela tem o caráter de língua oficial.
Trata -se, portanto, de definir a opção pelo ensino da língua \textbf{espanhola} como língua franca, uma língua de comunicação internacional utilizada por falantes espalhados no mundo inteiro, com diferentes repertórios linguísticos e culturais.
Essa perspectiva permite questionar a visão de que o único \textbf{espanhol} correto – e a ser ensinado – é aquele falado por \textbf{espanhóis}, por exemplo.
Desse modo, o tratamento do \textbf{espanhol} como língua franca o desvincula da noção de pertencimento a um determinado território e, consequentemente, há culturas típicas de comunidades específicas.
Esse entendimento favorece uma educação linguística voltada para a interculturalidade, isto é, para o reconhecimento das ( e o respeito às ) diferenças, e para a \textbf{compreensão} de como elas são produzidas.
A segunda implicação \textbf{diz} respeito à ampliação da visão de \textbf{letramento}, ou melhor, dos \textbf{letramentos}, concebida especialmente nas práticas sociais do mundo \textbf{digital} – no qual saber a língua \textbf{espanhola} potencializa as possibilidades de participação e circulação – que \textbf{aproximam} e entrelaçam diferentes semioses e linguagens ( verbal, visual, corporal, audiovisual ).
Essas práticas criam novas possibilidades de identificar e expressar ideias, sentimentos e valores.
O ensino de uma língua estrangeira deve partir não somente dos conteúdos linguísticos correspondentes a ela, mas, também, correlacionar o seu aprendizado aspectos vinculados à língua como um todo.
Esses aspectos correspondem a questões relativas ao povo que fala e vive o idioma em estudo, tais como : geográficos, culturais, \textbf{econômicos}, sociopolíticos e outros, tornando o processo aprendizagem mais rico e significativo para aquele que aprende.
Com relação à aprendizagem da língua \textbf{espanhola}, como segunda língua, em sala de \textbf{aula}, tem -se um campo vasto de conteúdos relacionados a ela, além dos linguísticos.
No caso específico, são instrumentos que com ela relacionam -se.
Um deles é o Mercosul ( Mercado Comum do Sul ), visto como elemento de integração dos países do Cone Sul.
A Constituição Federal de 1988 trata em seu \textbf{Parágrafo} Único, do Artigo 4o. : “ A República Federativa do Brasil buscará a integração \textbf{econômica}, política, social e cultural dos povos da América Latina, visando à formação de uma comunidade latino-americana de nações ”.
É de \textbf{suma} importância e constitucional esta integração entre os países e suas línguas, onde o Brasil e seus entes federados, não podem esquecer de incluir no sistema educativo a língua \textbf{espanhola}.
Outra relação que se deve ter no nosso sistema educacional a partir da língua \textbf{espanhola} é o acolhimento dos refugiados latino-americanos, haja \textbf{vista} que muitos que vêm ao nosso Estado do Piauí, são hispano-falantes.
Os esforços de apoio no nível regional devem ser redobrados em vários aspectos, mas especialmente no contexto da assistência humanitária e na cooperação internacional.
É fundamental proporcionar a inclusão social laboral dos refugiados e migrantes para que acessem o trabalho formal sem exploração laboral ( venezuelanos, bolivianos e peruanos, por exemplo ) sem uma queda de salários para os cidadãos nacionais das comunidades receptoras.
Um deles é a língua materna ( \textbf{espanhol} ) ajudando com a língua portuguesa.
De acordo com Santos ( 2017 ), professor da Universidade Estadual do Piauí, ex-presidente da Associação de Professores de \textbf{Espanhol} do Estado de Alagoas e \textbf{atual} presidente da Associação de Professores de \textbf{Espanhol} do Estado do Piauí, “ A aprendizagem da Língua \textbf{Espanhola} é uma possibilidade de aumentar a autopercepção do \textbf{aluno} como ser humano e como cidadão.
Por esse motivo, ela deve centrar -se no engajamento \textbf{discursivo} do \textbf{aprendiz}, ou seja, em sua \textbf{capacidade} de se engajar e engajar outros no discurso de modo a poder agir no mundo social ”.
Para que isso seja possível, é fundamental que o ensino de Língua Estrangeira, no caso do \textbf{espanhol}, seja balizado pela função social desse conhecimento na sociedade brasileira, e principalmente no Estado em tela.
Tal função está, principalmente, relacionada ao uso que se faz de Língua Estrangeira via leitura, embora se possa também considerar outras habilidades comunicativas em função da especificidade de algumas línguas estrangeiras e das condições existentes no contexto escolar.
Além disso, em uma política de pluralismo linguístico, condições pragmáticas apontam necessidade de considerar três fatores para orientar a inclusão de uma determinada língua estrangeira no currículo : fatores relativos à história, às comunidades locais e à tradição.
Duas questões \textbf{teóricas} \textbf{ancoram} os \textbf{parâmetros} de Língua Estrangeira : uma visão sociointeracional da linguagem e da aprendizagem.
O enfoque sociointeracional da linguagem indica que, ao se engajarem no discurso, as pessoas consideram aqueles a quem se dirigem ou quem se dirigiu a elas na construção social do significado.
É determinante nesse processo o \textbf{posicionamento} das pessoas na instituição, na cultura e na história.
Para que essa natureza sociointeracional seja possível, o \textbf{aprendiz} utiliza conhecimentos sistêmicos, de mundo e sobre a organização textual, além de ter de aprender como usá -los na construção social do significado via Língua Estrangeira.
A consciência desses conhecimentos e a de seus usos são essenciais na aprendizagem, posto que focaliza aspectos metacognitivos e desenvolve a consciência \textbf{crítica} do \textbf{aprendiz} no que se refere a como a linguagem é usada no mundo social, como reflexo de crenças, valores e projetos políticos.
No que se refere à visão sociointeracional da aprendizagem, pode -se \textbf{dizer} que é compreendida como uma forma de se estar no mundo com alguém e é, igualmente, situada na instituição, na cultura e na história.
Assim, os processos cognitivos têm uma natureza social, sendo gerados por meio da interação entre um \textbf{aluno} e um parceiro mais competente.
Em sala de \textbf{aula}, esta interação tem, em geral, caráter assimétrico, o que coloca dificuldades específicas para a construção do conhecimento.
Daí a importância de o professor aprender a compartilhar seu poder e dar \textbf{voz} ao \textbf{aluno} de modo que este possa se constituir como sujeito do discurso e, portanto, da aprendizagem.
Os temas \textbf{centrais} desta proposta são a cidadania, a consciência \textbf{crítica} em relação à linguagem e os aspectos sociopolíticos da aprendizagem de Língua Estrangeira.
Esses temas se articulam com os temas \textbf{transversais} dos \textbf{Parâmetros} Curriculares Nacionais, notadamente, na possibilidade de se usar a aprendizagem de línguas como espaço para se compreender, na escola, as várias \textbf{maneiras} de se viver a experiência humana O Componente Curricular de Língua \textbf{Espanhola} foi definido pela Base Nacional Comum Curricular ( BNCC – homologada em 2017 ) como o idioma estrangeiro optativo a ser ensinado para a Educação Básica a partir do Ensino Médio de todas as escolas brasileiras.
A língua \textbf{Espanhola} é integrante da área de linguagens, que tem em sua composição 07 \textbf{competências} e 28 habilidades no Ensino Médio.
Compreendida como língua de uso mundial, pela multiplicidade e variedade de usos, usuários e funções na \textbf{contemporaneidade}, o \textbf{espanhol} é legitimado pela Base como uma oportunidade de acesso ao mundo \textbf{globalizado}, deixando de ser apenas dos falantes nativos, o que favorece o ensino da língua com mais interculturalidade, porque varia com diferentes contextos.
Conforme BNCC ( 2017, p. 476, grifo nosso ), “ IX – língua \textbf{inglesa}, podendo ser oferecidas outras línguas estrangeiras, em caráter optativo, preferencialmente o \textbf{Espanhol}, de acordo com a disponibilidade da instituição ou rede de ensino ( Resolução CNE/CEB no 3/2018, Art. 11, § 4o ) ”.
O texto da Base define ainda como objetivo para o ensino de língua \textbf{espanhola} que os estudantes devem desenvolver \textbf{competências} e habilidades a partir de uma perspectiva de educação linguística consciente, \textbf{crítica} e \textbf{reflexiva}, de forma que a aprendizagem do idioma propicie aos estudantes o acesso a novos percursos de construção de conhecimento e o exercício da cidadania ativa, permitindo -lhes vivenciar “ novas formas de engajamento e participação em um mundo social cada vez mais \textbf{globalizado} e plural ” ( BNCC, 2017, p. 241 ).
Assim, o ensino de língua \textbf{espanhola} está \textbf{ancorado} na perspectiva de uma \textbf{abordagem} voltada para a oralidade, audição, leitura e escrita para o seu uso em um viés que transcende uma aprendizagem voltada para dimensão linguística, passando ao universo das dimensões sociais e culturais.
Considerando que o público do Ensino Médio é o jovem e que nesta fase afloram dúvidas naturais em relação à construção da identidade, além dos aspectos citados acima, o ensino de língua \textbf{espanhola} bem como dos demais componentes curriculares deve estar balizado na premissa de que o jovem tem características peculiares e, portanto precisa garantir o protagonismo, a \textbf{pluralidade} e diversidade próprias dessa faixa etária, proporcionando autonomia e formação de estudantes \textbf{críticos} com \textbf{capacidade} para definir seus projetos de vida.
De acordo com a BNCC ( 2017, p. 464 ), “ Para atender às necessidades de formação geral, \textbf{indispensáveis} ao exercício da cidadania e à inserção no mundo do trabalho, e responder à diversidade de expectativas dos jovens quanto à sua formação, a escola que acolhe as \textbf{juventudes} têm de estar comprometida com a educação integral dos estudantes e com a construção de seu projeto de vida ”.
Desta forma, este texto que apresenta a \textbf{contextualização} para o ensino de Língua \textbf{Espanhola} no Ensino Médio está sequenciado do contexto histórico da oferta e os marcos legais que fundamentam o ensino da língua no país e no estado do Piauí, apresentando uma \textbf{abordagem} \textbf{crítica} da inserção da língua no currículo, sendo considerada inicialmente como língua estrangeira \textbf{moderna} e, posteriormente assumindo o status de língua franca, pois permite a integração com grupos multilíngues e multiculturais no mundo global.
Portanto, \textbf{esperamos} que estudantes, professores e a comunidade escolar piauiense façam uso da língua em diferentes contextos interacionais e comunicativos, possibilitando o compartilhamento de experiências e abrindo \textbf{caminhos} para que o aprendizado e a aplicação da língua aconteçam de forma natural e fruitiva e ainda, que seja meio para ampliar o engajamento \textbf{discursivo} nas redes sociais e demais softwares que permitam as trocas de informação no mundo interativo e \textbf{globalizado}.
Em seguida são apresentados aspectos que constituem a aprendizagem da Língua \textbf{Espanhola} e a organização curricular com \textbf{abordagem} das \textbf{competências} e habilidades gerais e específicas da língua para o Ensino Médio.
O Desenvolvimento de \textbf{Competências} e Habilidades na Área de Linguagens O presente \textbf{tópico} aborda alguns aspectos necessários à \textbf{compreensão} do significado de \textbf{competências} e habilidades, no contexto da área de Linguagem e suas Tecnologias, no \textbf{intuito} de \textbf{auxiliar} o profissional de educação no trabalho com os estudantes.
As dez \textbf{competências} gerais sumarizam os direitos de aprendizagem e desenvolvimento dos estudantes, sendo fundamental que todas sejam manifestadas até o \textbf{final} da etapa do Ensino Médio.
No que tange às \textbf{competências} especificas da área de Linguagem e do Componente Curricular ( Língua Portuguesa ) é fundamental que sejam potencializadas de modo articulado com as \textbf{competências} gerais, que significa na prática, envolver as quatro dimensões da educação do \textbf{século} XXI : aprender a aprender, aprender a fazer, aprender a ser e aprender a conviver.
As \textbf{Competências} Específicas da Área de Linguagens guardam relação com conceitos diretamente relacionados aos campos dos saberes dos componentes curriculares, Língua Portuguesa, Língua estrangeira, Arte e Educação Física.
Nesse processo, serão \textbf{mobilizados} os conhecimentos essenciais de cada componente da área, \textbf{imprescindíveis} para a consolidação e ampliação das habilidades de uso e de reflexão sobre as linguagens – artísticas, corporais e verbais ( oral ou \textbf{visual-motora}, como \textbf{Libras}, e escrita ) que se constituem como objeto de estudos como prevê BNCC ( 2018 ).
Essas \textbf{competências} específicas da área em articulação com os campos de atuação buscam orientar as práticas de linguagem visando o desenvolvimento integral dos estudantes para que possam manifestar -se através de respostas inéditas, criativas, considerando as linguagens artísticas, corporais, linguísticas e \textbf{digitais} ; e que se traduzem em compromisso de atuação com as novas realidades.
As habilidades expressam as aprendizagens essenciais que devem ser asseguradas aos \textbf{alunos} nos diferentes contextos escolares e devem estar correlacionadas ao \textbf{escopo} de aprendizagem \textbf{esperados} pela \textbf{competência} específica.
O Objeto do conhecimento \textbf{diz} respeito à seleção de conteúdos que o professor pode lançar mão para desenvolver tais habilidades com os estudantes.
Objetivos da aprendizagem seria o \textbf{desdobramento} dessas habilidades pelo professor com a finalidade de movimentar o estudante no sentido de desenvolver uma aprendizagem autônoma com a mobilidade necessária para se relacionar com os conteúdos, apropriar -se e utilizá -los quando as diversas situações cotidianas \textbf{demandarem}.
O Documento Curricular de Referência ( DCR ) do Estado do Piauí apresenta a área de Linguagens composta por sete \textbf{Competências} Específicas da Área que se encontram articuladas às dez \textbf{Competências} Gerais ( conforme \textbf{quadro} a seguir ) às quais serão tratadas de forma transdisciplinar.
No Ensino Médio, a área tem a responsabilidade de propiciar oportunidades para a consolidação e ampliação das habilidades de uso e de reflexão sobre as linguagens – artísticas, corporais e verbais ( oral ou \textbf{visual-motora}, como \textbf{Libras}, e escrita ) –, que são objeto de seus diferentes componentes ( Arte, Educação Física, Língua \textbf{Inglesa} e Língua Portuguesa ( BNCC, p. 483 ).

% matched lemmas: abordagem, aluno, ancorar, aprendiz, aproximar, atual, aula, auxiliar, caminho, capacidade, central, competência, compreensão, contemporaneidade, contextualização, crítico, demandar, desdobramento, digital, discursivo, dizer, econômico, escopo, espanhol, esperar, final, globalizado, horizonte, imprescindível, indispensável, inglês, intuito, juventude, letramento, libra, maneira, mobilizar, moderno, parágrafo, parâmetro, pluralidade, posicionamento, quadro, reflexivo, sumo, século, teórico, transversal, tópico, vista, visual-motor, voz
\end{textsample}
